\chapter{بحث و گفتگو}
\label{ch:discussion}

%==============================================================================
\section{تفسیر عملکرد}
%==============================================================================

\subsection{از نتیجۀ منفی تا نزدیکی به تساوی}

یافتۀ کلیدی این پایان‌نامه این است که خط پایۀ ساده‌ی کم‌نمونه می‌تواند به‌طور محسوس از مدل‌های جدولی قوی روی همان وظیفه عقب بماند. این موضوع صریح گزارش شده است، زیرا نتایج منفی نیز مرزهای روش را روشن می‌کنند و از تکرار مسیرهای ناموفق جلوگیری می‌نمایند. بر اساس \textbf{معیارهای منجمد ارسال}، مدل جدولی \lr{SOTA} به \lr{AUROC} آزمون \textbf{\pnum{0.748}} می‌رسد، در حالی که اجرای انباشتی یکپارچه مقدار \textbf{\pnum{0.746}} \cnum{[0.699, 0.797]} را ثبت می‌کند؛ یعنی تنها حدود \textbf{\pnum{0.002}} پایین‌تر از برآورد نقطه‌ای \lr{SOTA}. بهترین زیرمدل کم‌نمونه به‌تنهایی \textbf{\pnum{0.670}} است. مقایسۀ \lr{Exp1} روی تقسیم ثابت نیز شکاف کنترل‌شده را نشان می‌دهد: مدل پیشنهادی کم‌نمونه با \textbf{\pnum{0.611}} در برابر \lr{RF} با \textbf{\pnum{0.762}}. فصل~\ref{ch:results} علاوه بر این، کمپین \lr{V1--V6} را به‌عنوان مسیر تاریخی توسعه مستند می‌کند.

در تفسیر بالینی، امکان‌پذیری درمان در این پایان‌نامه یک برچسب عملیاتی مبتنی بر پیامد ترخیص/بستری مجدد است، نه توصیه مستقیم درمان یا جایگزین قضاوت پزشک. بنابراین حتی بهترین مدل فقط می‌تواند بیمار را از نظر احتمال قرارگرفتن در مسیر قابل‌قبول درمان/ترخیص رتبه‌بندی کند و برای کاربرد واقعی به اعتبارسنجی محلی، کالیبراسیون و نظارت بالینی نیاز دارد.

\subsection{سقف تجربی و بستن تدریجی فاصله}

نتایج ما سقف جدولی روشنی برای این کوهورت نشان می‌دهد: \textbf{\lr{AUROC}} بسته به پروتکل در بازۀ \textbf{\prange{0.748}{0.760}} قرار دارد. مقدار \pnum{0.762} برای \lr{Random Forest} خروجی پروتکل ثابت \lr{Exp1} است، اما مقدار \pnum{0.748} مربوط به فرایند منجمد \lr{Tabular SOTA} است؛ بنابراین این دو عدد را نباید به‌عنوان تناقض خواند. پیکربندی \emph{تاریخی} آغازین در کمپین توسعه با کم‌نمونه ساده نزدیک \textbf{\pnum{0.619}} شروع می‌شود؛ اما این عدد با \lr{Exp1} ثابت (\pnum{0.611}) یکسان نیست و نباید جایگزین آن شود. راهبردهای بعدی در همان مسیر توسعه، با انباشت، \lr{GFA} و آموزش چندمرحله‌ای فاصله را کاهش می‌دهند. نتیجۀ انباشتی کم‌نمونه که مرجع ارسال است، \textbf{\pnum{0.746}} است.

%==============================================================================
\section{محدودیت‌های فرا یادگیری کم‌نمونه}
%==============================================================================

\subsection{ناسازگاری اپیزودیک}

امکان‌پذیری درمان ممکن است با فرا یادگیری اپیزودیک هم‌راستا نباشد. وظیفه ممکن است به آمار سراسری (مثلاً توزیع آزمایشگاهی جمعیت، عوامل ریسک سطح کوهورت) نیاز داشته باشد نه نمونه‌های اولیۀ محلی مجموعۀ پشتیبان. شبکه‌های نمونه‌اولیه فرض می‌کنند هویت کلاس را میانگین چند نهفته‌سازی پشتیبان می‌گیرد؛ برای امکان‌پذیری درمان، مرز تصمیم ممکن است به تعاملات غیرخطی پیچیده وابسته باشد که با مرکزهای نمونه‌اولیه خوب تقریب زده نمی‌شوند.

برای تصویرسازی، دو بیمار را در نظر بگیرید: بیمار \lr{A} کراتینین بالا (\pnum{0.8} نرمال‌شده) و لاکتات نرمال (\pnum{0.2}) دارد؛ بیمار \lr{B} کراتینین نرمال (\pnum{0.2}) و لاکتات بالا (\pnum{0.7}) دارد. هر دو ممکن است غیرممکن باشند اما به دلایل متفاوت. نمونه‌اولیۀ کلاس «غیرممکن» این پروفایل‌های متفاوت را میانگین می‌گیرد و مرکزی با کراتینین متوسط (\pnum{0.5}) و لاکتات متوسط (\pnum{0.45}) می‌سازد؛ پروفایلی که \emph{هیچ‌کدام} از بیماران را خوب نمایندگی نمی‌کند. این مشکل «میانگین‌گیری نمونه‌اولیه» بنیادین برای شبکه‌های نمونه‌اولیه است و وقتی کلاس‌ها توزیع چندوجهی یا ناهمگن دارند تشدید می‌شود. مدل‌های درخت‌بنیاد این وضعیت را طبیعی‌تر از راه افرازهای ویژگی‌محور که زیرجمعیت‌های مختلف را در یک کلاس جدا می‌کنند، مدیریت می‌کنند.

\subsection{ناسازگاری سوگیری القایی}

دادۀ جدولی و تصویر ساختار بنیاداً متفاوتی دارند و سوگیری‌های القایی که کم‌نمونه را در تصویر موفق می‌کند ممکن است به جدول منتقل نشود:

\begin{enumerate}
\item \textbf{همسایگی فضایی در تصاویر در برابر استقلال ویژگی در جدول.} شبکه‌های کانولوشنی از همسایگی فضایی؛ پیکسل‌های نزدیک همبسته؛ برای یادگیری سلسله‌مراتبی بهره می‌برند. ویژگی‌های جدولی این ساختار را ندارند؛ کراتینین و ضربان قلب در هیچ معنای فضایی «مجاور» نیستند. مکانیزم خودتوجه ترنسفورمر می‌تواند تعاملات ویژگی دلخواه یاد بگیرد اما باید از صفر بدون سوگیری فضایی این کار را کند.
\item \textbf{شکاف یادگیری انتقالی.} در بینش، روش‌های کم‌نمونه از ستون فقرات‌های ازپیش‌آموزش‌دیده (\lr{ImageNet}) سود می‌برند. ستون فقرات مشابهی برای دادۀ جدولی بالینی وجود ندارد. رمزگذار \lr{ETHOS} ما از صفر روی \pnum{1,280} بیمار آموزش دیده که کیفیت بازنمایی‌های آموخته‌شده را محدود می‌کند.
\item \textbf{مهندسی ویژگی به‌عنوان دانش ضمنی قبلی.} مدل‌های کلاسیک روی ویژگی‌های دست‌ساز مانند \lr{shock index} و پراکسی \lr{SOFA} کار می‌کنند؛ ویژگی‌هایی که دهه‌ها دانش بالینی را در خود دارند. رمزگذار \lr{ETHOS} باید این روابط را از دادۀ محدود دوباره کشف کند. شکاف میان \lr{ETHOS} تاریخی (\pnum{0.619}) و \lr{LR} روی ویژگی‌های خام (\pnum{0.724}) نشان می‌دهد فرایند کدگذاری، دست‌کم در پیکربندی‌های اولیه، ممکن است بیش از آنکه از تعاملات آموخته‌شده سود ببرد بخشی از اطلاعات را از دست بدهد.
\item \textbf{تنوع وظیفه در فرا آموزش.} بنچمارک‌های کم‌نمونه (\lr{miniImageNet}) از صدها کلاس نمونه می‌گیرند و اپیزودهای متنوع می‌سازند. طبقه‌بندی \pnum{2}راه با \pnum{5} گرۀ بیماری تنوع وظیفۀ محدودی دارد که ممکن است به فرا-بیش‌برازش منجر شود؛ مدل توزیع اپیزود خاص را حل می‌کند نه ترکیب‌های جدید بیمار را تعمیم می‌دهد.
\end{enumerate}

\subsection{مقیاس داده}

\pnum{2,000} بیمار ممکن است برای فرا یادگیری جهت یادگیری بازنمایی‌های تاب‌آور کافی نباشد. روش‌های کم‌نمونه اغلب به مجموعه‌های فرا آموزش بزرگ برای شبیه‌سازی وظایف متنوع نیاز دارند. با تنها \pnum{2,000} بیمار و \pnum{5} گرۀ بیماری، تنوع اپیزودها محدود است. پیش‌آموزش بزرگ‌تر روی \lr{MIMIC\text{-}IV} یا \lr{EHR} خارجی می‌تواند عملکرد را بهبود دهد.

رابطۀ بین اندازۀ مجموعۀ فرا آموزش و \lr{AUROC} مورد انتظار را می‌توان از تحلیل حساسیت $K$-شات و مقایسه با مدل‌های کلاسیک به‌صورت کیفی تفسیر کرد. مدل‌های کلاسیک در این کوهورت زودتر پایدار می‌شوند، در حالی که چارچوب \lr{ETHOS} و یادگیری کم‌نمونه واریانس بالاتر و عملکرد مطلق پایین‌تر نشان می‌دهد. بنابراین محتمل است که مسیر کم‌نمونه برای رسیدن به عملکرد مدل‌های جدولی قوی به پیش‌آموزش بزرگ‌تر، داده‌ی چندسایتی و اصلاح گلوگاه بازنمایی نیاز داشته باشد. مقدار دقیق داده‌ی لازم از نتایج حاضر قابل تعیین نیست و باید در آزمایش مقیاس‌پذیری مستقل سنجیده شود.

\subsection{بازنمایی}

توکن‌های علائم ممکن است نسبت به ویژگی‌های جدولی خام اطلاعات را از دست بدهند. مدل‌های کلاسیک بردار کامل \pnum{25}بعدی را مستقیماً دریافت می‌کنند؛ \lr{ETHOS} باید از دنباله‌های توکن اطلاعات مرتبط استخراج کند. سوگیری القایی ترنسفورمر ممکن است با ساختار وظایف پیش‌بینی بالینی هم‌تراز نباشد.

%==============================================================================
\section{تحلیل تفصیلی: چرا خط پایۀ اولیۀ کم‌نمونه ضعیف‌تر بود}
%==============================================================================

شکاف اولیه \pnum{0.14} \lr{AUROC} بین \lr{RF} (\pnum{0.760}) و خط پایۀ \lr{ETHOS} با یادگیری کم‌نمونه (\pnum{0.619}) نیازمند تحلیل دقیق است. درک این عوامل راهبرد بهبود ما را شکل داد. \textbf{پنج عامل زیر}، با تکیه بر شواهد حذف تدریجی:

\subsection{عامل \pnum{1}: گلوگاه فشرده‌سازی بازنمایی}

رمزگذار \lr{ETHOS} ورودی \pnum{25}بعدی را به نهفته‌سازی \pnum{128}بعدی فشرده می‌کند، سپس سر نمونه‌اولیه آن را به فاصلۀ اسکالر کاهش می‌دهد. با وجود آن‌که بعد نهفته (\pnum{128}) از بعد ورودی (\pnum{25}) بیشتر است، تبدیل بی‌زیان نیست. نتیجۀ حذف تدریجی: \lr{LR} روی \pnum{25} ویژگی خام \lr{AUROC} \pnum{0.724} دارد، بالاتر از \lr{ETHOS} روی نهفته‌سازی \pnum{128}بعدی (\pnum{0.619})؛ نشان می‌دهد رمزگذار \lr{ETHOS} صرفاً اطلاعات را حفظ نمی‌کند بلکه فعالانه \emph{تبدیل} می‌کند و این تبدیل ممکن است ویژگی‌های حیاتی برای مرز طبقه‌بندی را حذف کند.

از دید روش‌شناختی، این نتیجه به‌معنای آن نیست که نهفته‌سازی \lr{ETHOS} همیشه ضعیف‌تر از ویژگی خام است؛ بلکه نشان می‌دهد در این کوهورت و با این اندازۀ داده، تبدیل توکنی و سر نمونه‌اولیه‌ای هنوز نتوانسته‌اند تمام سیگنال جدولی لازم برای مرز تصمیم را حفظ کنند. بنابراین تفسیر درست، «گلوگاه بازنمایی» است، نه ناکارآمدی ذاتی معماری.

\subsection{عامل \pnum{2}: هندسۀ نمونه‌اولیه در فضاهای کم‌بعد}

شبکه‌های نمونه‌اولیه با نزدیک‌ترین نمونه‌اولیه در فضای نهفته طبقه‌بندی می‌کنند. برای کارایی خوب، فضای نهفته باید: (\pnum{1}) \emph{فشردگی درون‌کلاسی}: نمونه‌های هم‌کلاس محکم دور نمونه‌اولیه خوشه شوند؛ و (\pnum{2}) \emph{جداسازی بین‌کلاسی}: نمونه‌های اولیۀ کلاس‌های مختلف دور از هم باشند. تجسم‌های \lr{t-SNE} ما (فصل~\ref{ch:results}) جداسازی متوسط اما ناقص با هم‌پوشانی قابل‌توجه در ناحیۀ تصمیم نشان می‌دهند.

به‌صورت رسمی، نرخ خطای طبقه‌بندی نمونه‌اولیه توسط
\begin{equation}
P(\text{\lr{error}}) \leq \exp\left(-\frac{d(\mathbf{c}_0, \mathbf{c}_1)^2}{8\sigma^2}\right)
\end{equation}
کران داده می‌شود که در آن $d(\mathbf{c}_0, \mathbf{c}_1)$ فاصله‌ی بین نمونه‌های اولیه و $\sigma^2$ واریانس درون‌کلاسی است. وقتی فاصله‌ی نمونه‌های اولیه کوچک و واریانس درون‌کلاسی زیاد باشد، مرز تصمیم ناپایدار می‌شود؛ این دقیقاً همان رفتاری است که در نتایج اپیزودیک کم‌دامنه دیده می‌شود. \lr{RF} برعکس در فضای ویژگی اصلی کار می‌کند و با افرازهای هم‌محور می‌تواند زیرجمعیت‌های بالینی متفاوت را صریح‌تر جدا کند.

\subsection{عامل \pnum{3}: واریانس نمونه‌برداری اپیزودیک}

هر اپیزود آموزش $K$ پشتیبان و $Q$ پرس‌وجو به‌ازای هر کلاس نمونه می‌گیرد. با $K=5$، نمونه‌اولیه از \pnum{5} بیمار محاسبه می‌شود؛ برآورد پرنویز از مرکز کلاس. واریانس برآورد نمونه‌اولیه
\begin{equation}
\text{\lr{Var}}(\mathbf{c}_k) = \frac{\sigma_k^2}{K}
\end{equation}
است که $\sigma_k^2$ واریانس درون‌کلاسی است. با $K=5$، واریانس نمونه‌اولیه $\sigma_k^2/5$ است؛ \textbf{\pnum{5} برابر} بیش از میانگین جمعیت. این نویز به مرز طبقه‌بندی منتقل می‌شود و نرخ خطا را افزایش می‌دهد. در عمل، جاروب \lr{K-shot} ما فقط تغییرات کوچکی نشان می‌دهد: \lr{AUROC} اپیزودیک در بازۀ \prange{0.557}{0.564} می‌ماند و بهترین مقدار در \lr{K=3} دیده می‌شود. بنابراین انتظار نظریِ کاهش واریانس با \lr{K} بزرگ‌تر در این وظیفه وجود دارد، اما اثر تجربی آن در مقیاس کنونی ضعیف است.

\subsection{عامل \pnum{4}: مدل‌سازی تعامل ویژگی}

\lr{Random Forest} تعاملات ویژگی را طبیعی از راه افراز درخت می‌گیرد: افراز روی کراتینین و سپس سن، تعامل «کراتینین و سن» را می‌گیرد. \lr{RF} عمق \pnum{28} با \pnum{200} درخت می‌تواند تعاملات تا مرتبه \pnum{28} را مدل کند. ترنسفورمر \lr{ETHOS} تعاملات را از راه خودتوجه چندسر می‌گیرد، اما با تنها \pnum{2} لایه و \pnum{8} سر، ظرفیت برای تعاملات پیچیده ممکن است ناکافی باشد؛ به‌ویژه با دادۀ آموزش محدود.

تقویت گرادیان، مانند \lr{XGBoost} و \lr{LightGBM}، \lr{AUROC} نزدیک به \lr{RF} (\cnum{0.73-0.75}) به‌دست می‌آورد. این نتیجه تأیید می‌کند که مدل‌سازی تعاملات با درخت‌ها برای این وظیفه مؤثر است. مدل ترکیبی با \lr{AUROC} برابر \pnum{0.717} نشان می‌دهد \lr{ETHOS} برخی تعاملات مکمل را می‌گیرد که در ویژگی‌های مهندسی‌شده‌ی \pnum{25}بعدی کاملاً دیده نمی‌شوند؛ با این حال، این سیگنال به‌تنهایی برای بستن شکاف کافی نیست.

\subsection{عامل \pnum{5}: منظره بهینه‌سازی}

آموزش اپیزودیک یک فرا-هدف (زیان روی اپیزودها) را بهینه می‌کند، نه آنتروپی متقاطع استاندارد روی مجموعۀ آموزش. فرا-هدف پرنویزتر است (هر اپیزود برآورد تصادفی از فرا-زیان است) و ممکن است منظره بهینه‌سازی متفاوتی داشته باشد. با دادۀ فرا آموزش محدود (\pnum{1,280} بیمار $\div$ \pnum{5} به‌ازای هر کلاس در هر اپیزود $\approx$ \pnum{128} اپیزود ممکن پیش از تکرار)، مدل ممکن است منظره زیان را به‌اندازه کافی کاوش نکند. مدل‌های کلاسیک یک هدف واحد واضح روی کل آموزش بهینه می‌کنند و راه‌حل‌های پایدارتر می‌دهند.

%==============================================================================
\section{معامله‌ها و پیامدهای عملی}
%==============================================================================

\subsection{\texorpdfstring{چه‌وقت از \lr{ETHOS} و یادگیری کم‌نمونه استفاده کنیم}{چه‌وقت از ETHOS و یادگیری کم‌نمونه استفاده کنیم}}

در تنظیمات فدرال و کم‌داده که حریم خصوصی حیاتی است و مدل‌های کلاسیک نمی‌توانند به‌صورت متمرکز آموزش ببینند (مثلاً همکاری بین‌مؤسسه با حاکمیت داده سخت)، رویکرد ما خط پایۀ عملی ارائه می‌کند. شکاف فدرال در آزمایش گزارش‌شده حدود \pnum{0.005} واحد \lr{AUROC} است؛ بنابراین در این شبیه‌سازی، حریم خصوصی با افت عملکردی محدود به دست آمده است.

\subsection{چه‌وقت از مدل‌های کلاسیک استفاده کنیم}

برای بیشینه دقت روی این وظیفه، \lr{Random Forest} یا \lr{XGBoost} تنظیم‌شده ترجیح داده می‌شود. توصیه می‌کنیم پژوهشگران هر دو را گزارش کنند: مدل‌های کلاسیک به‌عنوان مرز بالا، و چارچوب \lr{ETHOS} با یادگیری کم‌نمونه به‌عنوان خط پایۀ فدرال/کم‌داده. انتخاب به بستر استقرار (متمرکز در برابر فدرال، دسترسی به داده) بستگی دارد. جدول~\ref{tab:model_selection} راهنمای تصمیم ارائه می‌کند.

\begin{table}[htbp]
\centering
\caption{راهنمای انتخاب مدل}
\begin{tabularx}{\linewidth}{@{}R{2.9cm}R{4.1cm}Y@{}}
\toprule
سناریو & مدل توصیه‌شده & استدلال \\
\midrule
بیشینه دقت، متمرکز & جدولی \lr{SOTA} یا \lr{RF} تنظیم آستانه & \lr{AUROC} در حدود \prange{0.748}{0.760} (فصل~\ref{ch:results}) \\
بهترین ترکیب کم‌نمونه + کلاسیک & \lr{Stacked BEST} & \lr{AUROC} \textbf{\pnum{0.746}} \cnum{[0.699, 0.797]} \\
بهترین عصبی تکی (کمپین \lr{V}) & \lr{V5} \lr{GFA}+\lr{CrossAttn} & بهترین بذر \pnum{0.737}؛ میانگین \lr{Ensemble} \pnum{0.722} (تاریخی) \\
فدرال، حریم بحرانی & چیدمان فدرال (فصل~\ref{ch:results}) & شکاف حدود \pnum{0.005} \lr{AUROC} نسبت به متمرکز \\
گرۀ جدید کم‌داده & \lr{Proto-MAML} \lr{FewShot} & سازگاری زمان آزمون با پشتیبان کم \\
استقرار تولید & \lr{Stacked BEST} + کالیبراسیون + \lr{SHAP} & نمونۀ پژوهشی؛ اعتبارسنجی محلی \\
\bottomrule
\end{tabularx}
\label{tab:model_selection}
\end{table}

\subsection{از ترکیبی به انباشت: مسیر تا تساوی}

مدل \texttt{\lr{ETHOS\_RF\_Hybrid}} با \lr{AUROC} \pnum{0.717} در کمپین \lr{V} نشان داد ترکیب بردارهای نهفتۀ \lr{ETHOS} با ویژگی‌های جدولی می‌تواند فاصله تا خطوط پایۀ کلاسیک را کم کند. انباشت \lr{V4} با \pnum{0.732} و بهترین بذر \lr{V5} با \pnum{0.737} نیز نقش هم‌آمیزی و معماری را بیشتر بررسی کردند. اجرای انباشتی \textbf{منجمد ارسال}، یعنی \texttt{\lr{exp\_fewshot\_best}} با \lr{AUROC} \textbf{\pnum{0.746}}، روایت کم‌نمونه + کلاسیک را با برآورد نقطه‌ای \lr{SOTA} جدولی (\textbf{\pnum{0.748}}) هم‌تراز می‌کند.

%==============================================================================
\section{اعتبارسنجی بالینی و تحلیل علی: چارچوب پیشنهادی}
%==============================================================================

نتایج این پایان‌نامه ماهیت پیش‌بینی دارند، نه علی. بنابراین از آن‌ها نباید نتیجه گرفت که خود مدل باعث کاهش بستری مجدد، کاهش طول اقامت یا بهبود پیامد بالینی می‌شود. مدل تنها نشان می‌دهد با استفاده از داده‌های \pnum{24} ساعت نخست می‌توان بیماران را از نظر امکان‌پذیری درمان و ترخیص بهتر رتبه‌بندی کرد. برای تبدیل این یافته به ادعای بالینی، مرحلۀ بعدی باید یک اعتبارسنجی آینده‌نگر یا نیمه‌آزمایشی باشد.

\subsection{مسیر پیشنهادی برای اعتبارسنجی آینده}

برای ارزیابی اثر واقعی در محیط بالینی، سه گام پیشنهاد می‌شود: نخست، خروجی مدل روی کوهورت مستقل و بدون دخالت در تصمیم پزشک ثبت شود تا کالیبراسیون و عملکرد زیرگروهی بررسی گردد. دوم، بیماران با ریسک مشابه از نظر سن، جنس، تشخیص اصلی، شدت بیماری و شاخص‌های آزمایشگاهی با روش‌هایی مانند تطبیق امتیاز گرایش مقایسه شوند. سوم، اگر بیمارستان چند بخش یا چند مرکز با سیاست‌های متفاوت داشته باشد، می‌توان از طرح‌های نیمه‌آزمایشی مانند اختلاف در اختلاف‌ها یا متغیر ابزاری استفاده کرد؛ البته فقط زمانی که فرض‌های شناسایی به‌روشنی برقرار باشند.

\subsection{مرزبندی ادعا}

در نسخۀ حاضر، ادعای اصلی این است که مدل‌ها \emph{توان پیش‌بینی} امکان‌پذیری را دارند و در حالت انباشتی به سقف جدولی نزدیک می‌شوند. ادعای علی یا اثر مستقیم بر گردش‌کار درمانی گزارش نمی‌شود. این مرزبندی برای داوری علمی مهم است، زیرا حتی مدل دقیق نیز بدون مطالعۀ مداخله‌ای نمی‌تواند به‌تنهایی اثبات کند که تصمیم‌های درمانی را بهتر می‌کند.

%==============================================================================
\section{یادگیری چندوظیفه‌ای: مسیر توسعه}
%==============================================================================

یک مسیر طبیعی برای ادامۀ این پژوهش، گسترش رمزگذار \lr{ETHOS} به یادگیری چندوظیفه‌ای است؛ برای نمونه، امکان‌پذیری درمان، مرگ‌ومیر کوتاه‌مدت، طول اقامت \lr{ICU} و خطر بستری مجدد می‌توانند سرهای خروجی جداگانه داشته باشند و در یک رمزگذار مشترک آموزش ببینند. انتظار روش‌شناختی این است که وظایف نزدیک به هم، بخشی از سیگنال شدت بیماری و مسیر بهبود را مشترک کنند و در نتیجه بازنمایی پایدارتری بسازند.

با این حال، در این نسخه عدد چندوظیفه‌ای به‌عنوان نتیجۀ اصلی گزارش نمی‌شود؛ زیرا اجرای کامل آن نیازمند تعریف دقیق برچسب‌های ثانویه، توازن زیان‌ها، کنترل نشتی زمانی و اعتبارسنجی جداگانه برای هر پیامد است. بنابراین جایگاه چندوظیفه‌ای در این پایان‌نامه «پیشنهاد توسعه» است، نه ادعای عملکرد.

%==============================================================================
\section{تحلیل حساسیت}
%==============================================================================

با پیروی از چارچوب تحلیل نظری پیشنهاد، تحلیل‌های حساسیت برای تاب‌آوری یافته‌ها انجام شد.

\subsection{حساسیت ابرپارامتر}

ابرپارامترهای کلیدی را تغییر دادیم و تغییر \lr{AUROC} را اندازه گرفتیم. برای \lr{ETHOS}: بعد نهفته (\pnum{64}، \pnum{128}، \pnum{256})، تعداد سرها (\pnum{2}، \pnum{4}، \pnum{8})، و \lr{Dropout} (\pnum{0.1}، \pnum{0.2}، \pnum{0.3}). \lr{AUROC} در حدود \pnum{0.02} بالا یا پایین می‌شد و ترکیب \pnum{128} بعد و \pnum{4} سر، بهترین توازن را نشان داد. برای \lr{Random Forest}: \texttt{\lr{max\_depth}} (\cnum{5-25}) و \texttt{\lr{n\_estimators}} (\cnum{100-500}) بررسی شد. \lr{AUROC} از عمق \pnum{10} به بالا و با دست‌کم \pnum{200} برآوردگر، در بازۀ حدود \pnum{0.01} پایدار ماند.

\subsection{حساسیت شات پشتیبان \lr{K}}

با تغییر مقدار \lr{K} در میان \pnum{1}، \pnum{3}، \pnum{5}، \pnum{10} و \pnum{15} در آزمایش حساسیت \lr{K}-شات، \lr{AUROC} در بازۀ \prange{0.557}{0.564} باقی ماند. بهترین میانگین در \lr{K=3} با \pnum{0.564} دیده شد و \lr{K=1} تقریباً همان عملکرد را نشان داد. روند صعودی یکنواختی مشاهده نمی‌شود. مدل منجمد \texttt{\lr{Proposed\_FewShot}} در \lr{Exp1} روی ارزیابی ثابت عدد \pnum{0.611} دارد؛ بنابراین اختلاف عمده بیشتر از تفاوت پروتکل ارزیابی ناشی می‌شود تا صرفاً شمار شات پشتیبان.

\subsection{اثر ترکیب ماژول‌ها}

سهم هر ماژول را از راه حذف تدریجی \lr{leave-one-out} (جدول~\ref{tab:exp4}) تحلیل کردیم. حذف وزن‌دهی شدت \lr{AUROC} را به \pnum{0.577} کاهش داد، در حالی که حذف واگرایی زمانی عملاً تغییری ایجاد نکرد و روی \pnum{0.583} باقی ماند. به‌ویژه \textbf{\texttt{\lr{no\_symptom\_tokens}}}، یعنی \lr{Logistic Regression} روی \pnum{25} بعد خام، به \lr{AUROC} \pnum{0.724} رسید؛ بالاتر از مدل کامل \lr{ETHOS} با یادگیری کم‌نمونه با \pnum{0.583}. این نتیجه نشان می‌دهد بازنمایی مبتنی بر توکن، در مقایسه با ویژگی‌های جدولی خام، ممکن است برای این وظیفه بخشی از اطلاعات را از دست بدهد و رمزگذار \lr{ETHOS} باید این زیان را جبران کند. همچنین حذف سر کم‌نمونه و جایگزینی آن با سر نظارت‌شده استاندارد به \lr{AUROC} \pnum{0.645} می‌رسد که از مدل کامل بهتر است؛ بنابراین چیدمان اپیزودیک در این مقیاس الزاماً سودمند نیست.

%==============================================================================
\section{\texorpdfstring{اثر اقتصادی و بازگشت سرمایه‌ی توضیحی (\lr{ROI})}{اثر اقتصادی و بازگشت سرمایه توضیحی (ROI)}}
%==============================================================================

تحلیل اقتصادی در این پایان‌نامه به‌صورت چارچوب تصمیم‌گیری مطرح می‌شود، نه محاسبۀ قطعی بازگشت سرمایه. \lr{ROI} واقعی به ظرفیت \lr{ICU}، هزینه‌ی بستری مجدد، درصد پذیرش‌های قابل مداخله، هزینه‌ی اتصال به \lr{EHR}، بار کاری تیم درمان و سیاست‌های هر مرکز وابسته است. بنابراین به‌جای گزارش یک عدد ثابت، پیشنهاد می‌شود هر مرکز مقدار
\[
\mathrm{ROI}=\frac{\text{صرفه‌جویی ناشی از کاهش خطاهای ترخیص و مدیریت بهتر تخت}-\text{هزینۀ استقرار و پایش}}{\text{هزینۀ استقرار و پایش}}
\]
را با داده‌ی محلی خود برآورد کند. نقش نتایج این پژوهش در این محاسبه، فراهم کردن یک مدل رتبه‌بندی و چارچوب پایش است؛ اثبات اثر اقتصادی نیازمند مطالعۀ پیاده‌سازی واقعی است.

%==============================================================================
\section{گسترش سری زمانی: مسیر پیشنهادی}
%==============================================================================

توکن‌سازی \pnum{25}توکنی علامت، داده‌های \pnum{24} ساعت نخست را به بازنمایی ثابت و قابل تفسیر تبدیل می‌کند؛ اما بخشی از مسیر زمانی بیمار، مانند روند کاهش کراتینین یا پایدار شدن علائم حیاتی، ممکن است در این خلاصه‌سازی از دست برود. از این رو، مسیر بعدی می‌تواند افزودن مدل‌های سری زمانی مانند \lr{GRU}، \lr{LSTM} یا ترنسفورمر زمانی روی پنجره‌های چندساعته باشد.

در نسخۀ حاضر، مقایسۀ کامل و تنظیم‌شدۀ سری زمانی انجام نشده است؛ بنابراین هیچ برتری عددی برای \lr{LSTM} یا مدل زمانی دیگر گزارش نمی‌شود. پیام اصلی این بخش آن است که بازنمایی توکنی از نظر تفسیرپذیری، هزینۀ محاسباتی و سازگاری با کم‌نمونه/\lr{FL} مناسب است، اما برای بیشینۀ عملکرد آینده باید در کنار ویژگی‌های روندی و مدل‌سازی زمانی ارزیابی شود.

%==============================================================================
\section{تحلیل خطا}
%==============================================================================

جایی که بهترین مدل (جنگل تصادفی) خطا می‌کند را تحلیل کردیم. طبقه‌بندی اشتباه شایع‌تر میان بیماران با کراتینین مرزی (\prange{0.3}{0.5} نرمال‌شده) و در گره‌های بیماری با نمونۀ آموزش کمتر (کلیه، دیابت) است. کالیبراسیون (شکل~\ref{fig:calibration}) نشان می‌دهد احتمالات پیش‌بینی‌شده برای \lr{RF} نسبتاً خوب کالیبره‌اند؛ \lr{ETHOS} گرایش به بیش‌اعتمادی دارد. برای استقرار، نمودارهای قابلیت اطمینان و بازکالیبراسیون (مثلاً \lr{Platt scaling}) در صورت نیاز توصیه می‌شود.

%==============================================================================
\section{ملاحظات استقرار}
%==============================================================================

برای استقرار تولید توصیه می‌کنیم: (\pnum{1}) \textbf{یکپارچه‌سازی}: \lr{REST API} یا \lr{HL7 FHIR} برای \lr{EHR}؛ (\pnum{2}) \textbf{تأخیر}: استنتاج \mbox{\lr{RF $<{}$10\,ms}}؛ \mbox{\lr{ETHOS $<{}$50\,ms}} روی \lr{CPU}؛ (\pnum{3}) \textbf{پایش}: ردیابی \lr{AUROC}، کالیبراسیون، و معیارهای انصاف در زمان؛ (\pnum{4}) \textbf{پشتیبان}: بازبینی انسانی برای پیش‌بینی‌های کم‌اطمینان (با احتمال کمتر از \pnum{0.3} یا بیش از \pnum{0.7})؛ (\pnum{5}) \textbf{نسخه‌گذاری}: ثبت نسخۀ مدل و توزیع داده برای هر پیش‌بینی. مشخصات \lr{API} در پیوست~\ref{app:api} آمده است.

%==============================================================================
\section{\texorpdfstring{مشخصات رابط برنامه‌نویسی کاربردی (\lr{API})}{مشخصات رابط برنامه‌نویسی کاربردی (API)}}
%==============================================================================

برای استقرار و بازتولیدپذیری، مشخصات \lr{API} (پیوست~\ref{app:api}) ارائه شده است. \lr{REST API} نقاط پایانی برای: (\pnum{1}) پیش‌بینی امکان‌پذیری (\lr{POST} \texttt{\lr{/predict}})، (\pnum{2}) فراداده مدل (\lr{GET} \texttt{\lr{/model}})، و (\pnum{3}) توضیح \lr{SHAP} (\lr{POST} \texttt{\lr{/explain}}) دارد. برای مشخصات کامل \lr{OpenAPI} به پیوست~\ref{app:api} مراجعه کنید.

%==============================================================================
\section{تحلیل تفصیلی خطا به‌تفکیک زیرگروه}
%==============================================================================

نرخ خطا را به‌تفکیک گروه سنی (\prange{18}{45}، \prange{46}{65}، \pnum{66}+)، جنس و نوع پذیرش (انتخابی در برابر اورژانسی) تحلیل کردیم. برای بهترین مدل (\lr{RF})، نرخ خطا در گروه \pnum{66}+ بالاتر است (\pnum{24}\% در برابر \pnum{18}\% برای \prange{18}{45})، که با پیچیدگی بالینی بیشتر در سالمندان سازگار است. نرخ خطا به‌تفکیک جنس مشابه است (\pnum{20}\% در برابر \pnum{21}\%). پذیرش‌های اورژانسی نیز نرخ خطای بالاتری دارند (\pnum{23}\% در برابر \pnum{17}\% انتخابی)، که نشان می‌دهد شدت بیماری و ناپایداری مسیر بالینی می‌تواند پیش‌بینی را دشوارتر کند. این تحلیل‌های زیرگروهی برای ارزیابی انصاف مهم‌اند و در استقرار باید عملکرد مدل در گروه‌های جمعیت‌شناختی پایش شود.

%==============================================================================
\section{مقایسه با کار پیشین}
%==============================================================================

\subsection{ادبیات پیش‌بینی ترخیص}

کار پیشین روی پیش‌بینی مقصد ترخیص~\cite{awad2017random,caruana2015intelligible} معمولاً از دادۀ متمرکز و \lr{ML} کلاسیک استفاده می‌کند. \lr{AUROC} \pnum{0.76} ما با نتایج گزارش‌شده برای وظایف مشابه قابل‌قیاس است. جدول~\ref{tab:discharge_comparison} مقایسۀ تفصیلی ارائه می‌کند.

\begin{table}[htbp]
\centering
\caption{مقایسه با نتایج پیشین ترخیص/پیامد}
\footnotesize
\setlength{\tabcolsep}{2pt}
\renewcommand{\arraystretch}{1.16}
\begin{tabularx}{\linewidth}{@{}L{2.55cm}R{1.85cm}C{2.15cm}C{1.35cm}Y@{}}
\toprule
\textbf{مطالعه} & \textbf{وظیفه} & \textbf{مجموعه داده} & \textbf{\lr{AUROC}} & \textbf{روش کلیدی} \\
\midrule
\lr{Awad 2017} & مقصد ترخیص & \lr{MIMIC\text{-}III} & \cnum{0.78} & \lr{RF}، متمرکز \\
\lr{Purushotham 2018} & مرگ‌ومیر & \lr{MIMIC\text{-}III} & \cnum{0.85} & \lr{LSTM}، متمرکز \\
\lr{Rajkomar 2018} & بستری مجدد & چندسایتی & \cnum{0.76} & \lr{DNN}، متمرکز \\
\lr{Caruana 2015} & مرگ‌ومیر & \lr{EHR} & \cnum{0.88} & \lr{Ensembles} \\
\lr{Harutyunyan 2019} & چندوظیفه & \lr{MIMIC\text{-}III} & \cnum{0.86} & \lr{LSTM} چندوظیفه \\
\textbf{ما (\lr{RF}، \lr{SOTA})} & امکان‌پذیری & \lr{MIMIC\text{-}IV} & \cnum{0.748} & \lr{RF}؛ اسکریپت منجمد \lr{SOTA} \\
\textbf{ما (\lr{V5} بهترین)} & امکان‌پذیری & \lr{MIMIC\text{-}IV} & \cnum{0.737} & \lr{GFA}+\lr{CrossAttn} (کمپین \lr{V}) \\
\textbf{ما (انباشت \lr{V4})} & امکان‌پذیری & \lr{MIMIC\text{-}IV} & \cnum{0.732} & \lr{FS}+\lr{RF} (کمپین \lr{V}) \\
\textbf{ما (\lr{Stacked BEST})} & امکان‌پذیری & \lr{MIMIC\text{-}IV} & \textbf{\cnum{0.746}} & انباشت کم‌نمونه + کلاسیک \\
\textbf{ما (\lr{V6})} & امکان‌پذیری & \lr{MIMIC\text{-}IV} & \cnum{0.725} & \lr{Ensemble} ناهمگن (کمپین \lr{V}) \\
\textbf{ما (خط پایه منجمد)} & امکان‌پذیری & \lr{MIMIC\text{-}IV} & \cnum{0.611} & مدل پیشنهادی کم‌نمونه (\lr{Exp1}) \\
\bottomrule
\end{tabularx}
\label{tab:discharge_comparison}
\end{table}

نتیجۀ \lr{RF}/\lr{SOTA} در این پژوهش، بسته به پروتکل، در بازۀ \prange{0.748}{0.760} قرار می‌گیرد و با مقادیر گزارش‌شده برای وظایف مرتبط با ترخیص در دادۀ \lr{MIMIC} هم‌خوان است. نوآوری این کار در ادعای عبور از همۀ بنچمارک‌های پیشین نیست؛ بلکه در نشان دادن ظرفیت و محدودیت رویکردهای کم‌نمونه و فدرال برای یک مسئلۀ جدولی بالینی واقعی است.

\subsection{یادگیری کم‌نمونه در سلامت}

یادگیری کم‌نمونه در تصویربرداری پزشکی، مانند طبقه‌بندی \lr{X-ray} با \mbox{\lr{AUROC $>$ 0.85}} و پوست‌شناسی با دقت بیش از \mbox{\lr{80\%}}~\cite{prabhu2019few}، کاربردهای موفقی داشته است؛ اما در پیامدهای جدولی بالینی هنوز کمتر بررسی شده است. نقطۀ مشترک بسیاری از موفقیت‌های تصویری، وجود ستون‌فقرات‌های کانولوشنی ازپیش‌آموزش‌دیده است که پیش از سر کم‌نمونه، ویژگی‌های غنی استخراج می‌کنند. نتیجۀ منفی اولیه در این پژوهش، یعنی \lr{ETHOS} تاریخی با \pnum{0.619} در برابر \lr{RF} حدود \pnum{0.760}، نشان می‌دهد یادگیری کم‌نمونه جدولی احتمالاً یا به بازنمایی‌های ازپیش‌آموزش‌دیده‌ی قوی‌تر نیاز دارد یا به معماری‌هایی که مخصوص داده‌های جدولی ناهمگن طراحی شده باشند.

جهت‌های امیدبخش اخیر شامل:
\begin{enumerate}
\item \textbf{پیش‌آموزش \lr{EHR} مقیاس بزرگ}: پیش‌آموزش روی میلیون‌ها پرونده پیش از تنظیم دقیق کم‌نمونه، مشابه پیش‌آموزش \lr{ImageNet} برای بینش.
\item \textbf{پل‌های مدل زبانی}: استفاده از \lr{LLM} برای تولید توصیف متنی ویژگی‌های بیمار، سپس اعمال روش‌های کم‌نمونه مبتنی بر متن (\lr{TabLLM}~\cite{hegselmann2023tabllm}).
\item \textbf{نمونه‌های اولیه وابسته به وظیفه}: یادگیری توابع محاسبۀ نمونه‌اولیه ویژۀ وظیفه به‌جای میانگین ساده~\cite{ye2020few}.
\end{enumerate}

\subsection{بهداشت فدرال}

یادگیری فدرال در سلامت~\cite{rieke2020future,duan2020federated} در چند بستر پژوهشی و عملی امکان‌پذیر نشان داده شده است. شکاف کوچک فدرال در این پژوهش، حدود \pnum{0.005} واحد \lr{AUROC}، با گزارش‌های پیشین سازگار است؛ به‌ویژه زمانی که ناهمگنی داده متوسط باشد. برای نمونه، همکاری \lr{Intel--Penn Medicine}~\cite{sheller2020federated} گزارش کرد که بخش‌بندی فدرال تومور مغز در \pnum{30} مؤسسه به \pnum{99}\% عملکرد آموزش متمرکز رسیده است؛ نتیجه‌ای که امکان‌پذیری عملی \lr{FL} در سلامت را تقویت می‌کند.

تفاوت‌های کلیدی بین چیدمان فدرال شبیه‌سازی‌شده ما و \lr{FL} دنیای واقعی:
\begin{enumerate}
\item \textbf{قابلیت اطمینان ارتباط}: شبیه‌سازی ما ارتباط همزمان قابل‌اعتماد فرض می‌کند. \lr{FL} واقعی باید خرابی شبکه، \lr{timeout}، و به‌روزرسانی‌های ناهمزمان را مدیریت کند.
\item \textbf{ناهمگنی گره}: گره‌های ما بر اساس بیماری از هم جدا شده‌اند؛ اما در محیط واقعی، مؤسسات از نظر سخت‌افزار، نسخۀ نرم‌افزار و روند پردازش داده نیز با هم تفاوت دارند.
\item \textbf{شرکت‌کنندگان خصمانه}: گره‌های مخرب یا بیزانسی می‌توانند مدل سراسری را فاسد کنند. روش‌های تجمیع تاب‌آور (مثلاً میانگین هرس‌شده، \lr{Krum}~\cite{blanchard2017machine}) این را پوشش می‌دهند اما در خط پایۀ ما پیاده نشده‌اند.
\item \textbf{عقب‌ماندگان}: گره‌ها با سخت‌افزار کند یا دادۀ محلی بزرگ ممکن است دورها را به تعویق اندازند. روش‌های \lr{FL} ناهمزمان (\lr{FedAsync}) می‌توانند این را تخفیف دهند.
\end{enumerate}

%==============================================================================
\section{ملاحظات نظری}
%==============================================================================

\subsection{چرا شبکه‌های نمونه‌اولیه ممکن است ضعیف‌تر عمل کنند}

شبکه‌های نمونه‌اولیه فرض می‌کنند هویت کلاس را میانگین بردارهای نهفتۀ پشتیبان می‌گیرد. برای امکان‌پذیری درمان، مرز تصمیم ممکن است به تعاملات غیرخطی (مثلاً ترکیب کراتینین و سن، یا ترکیب \lr{SOFA} و تشخیص) وابسته باشد که با مرکزهای نمونه‌اولیه خوب تقریب زده نمی‌شوند. یادگیری متریک ممکن است به توابع فاصله ویژۀ وظیفه نیاز داشته باشد.

\subsection{کارایی داده}

با \pnum{2,000} بیمار و \pnum{5} گرۀ بیماری، حدود \pnum{400} بیمار به‌ازای هر گره داریم. فرا یادگیری کم‌نمونه اغلب از هزاران وظیفه برای نمونه‌برداری اپیزود متنوع سود می‌برد. مقیاس‌دهی به کل کوهورت \lr{MIMIC\text{-}IV} (\pnum{250}هزار+ اقامت) یا دادۀ چندسایتی می‌تواند عملکرد \lr{ETHOS} را بهبود دهد.

\subsection{تهدیدهای اعتبار}

چند تهدید اعتبار باید در نظر گرفته شود. از نظر \textbf{اعتبار درونی}، کوهورت از یک مؤسسه (\lr{BIDMC}) آمده و ممکن است ویژگی‌های همان مرکز را بازتاب دهد. از نظر \textbf{اعتبار بیرونی}، نمونۀ \pnum{2,000} بیمار نسبتاً محدود است و استفاده از کل \lr{MIMIC\text{-}IV} یا دادۀ چندسایتی می‌تواند به نتیجه‌گیری متفاوتی منجر شود. از نظر \textbf{اعتبار سازه}، امکان‌پذیری درمان با مقصد ترخیص و بستری مجدد تعریف شده است؛ تعاریف جایگزین مانند وضعیت عملکردی یا زمان تا ترخیص ممکن است نتایج را تغییر دهند. از نظر \textbf{اعتبار نتیجه‌گیری} نیز، چون همۀ مدل‌ها بردارهای پیش‌بینی هم‌تراز در یک پروتکل واحد تولید نمی‌کنند، استنتاج آماری جفتی برای همۀ مقایسه‌ها محدود است. برای کاهش این خطر، بازه‌های اطمینان موجود گزارش شده، از گزارش گزینشی پرهیز شده و مقایسه‌های بین‌پروتکلی به‌صورت توصیفی تفسیر شده‌اند.

\subsection{گلوگاه بازنمایی: تحلیل عمیق‌تر}

توکن‌سازی علائم، فضای جدولی \pnum{25}بعدی را به دنباله‌ای از توکن‌های گسسته تبدیل می‌کند. در این تبدیل، بخشی از اطلاعات ممکن است هنگام سطل‌بندی و تجمیع از دست برود. اینکه \lr{LR} روی ویژگی‌های خام با \lr{AUROC} برابر \pnum{0.724} از \lr{ETHOS} کامل با \pnum{0.583} بهتر عمل می‌کند، این نگرانی را تقویت می‌کند. این گلوگاه از سه منظر بررسی می‌شود:

\paragraph{منظر اطلاع‌نظری.} فرایند توکن‌سازی مقادیر پیوسته $x \in [0,1]$ را به شناسۀ علامت گسسته (\prange{1}{22}) و مقدار شدت پیوسته نگاشت می‌کند. با وجود حفظ مقدار نرمال‌شده توسط شدت، گسسته‌سازی هویت علامت واژگان ثابتی معرفی می‌کند که ممکن است تمایزهای ریز را نگیرد. برای مثال، مقادیر کراتینین \pnum{0.31} و \pnum{0.69} همان شناسۀ علامت را دارند اما پیامدهای بالینی متفاوت (نرمال در برابر بالا). مدل باید این تمایزها را تنها از بعد شدت یاد بگیرد، در حالی که مدل‌های کلاسیک مستقیماً روی مقدار پیوسته کار می‌کنند.

\paragraph{منظر معماری.} مکانیزم خودتوجه ترنسفورمر پیچیدگی محاسباتی $\mathcal{O}(L^2 d)$ دارد که در اینجا $L=25$ طول دنباله است. وقتی دنباله فقط \pnum{25} توکن دارد، پنجرۀ زمینه برای توجه محدود است و ممکن است نسبت به معماری‌های ساده‌تر، مانند \lr{MLP} دولایه روی بردار \pnum{25}بعدی صاف‌شده، مزیت روشنی ایجاد نکند. نتایج حذف تدریجی نیز همین نکته را تأیید می‌کند: جایگزینی ترنسفورمر با یک رمزگذار \lr{MLP} ساده \lr{AUROC} را تنها حدود \pnum{0.02} از \lr{ETHOS} کامل دور می‌کند. بنابراین، برای این طول دنباله و این وظیفه، خودتوجه به‌تنهایی عامل اصلی عملکرد نیست.

\paragraph{راهبردهای تخفیف.} بازنمایی‌های جایگزین که می‌توانند گلوگاه را پوشش دهند:
\begin{enumerate}
\item \textbf{بردارهای نهفتۀ توکن پیوسته}: جایگزینی شناسه‌های علامت گسسته با بردارهای نهفتۀ پیوسته آموزش‌دیده به‌ازای هر اندازه‌گیری، تا مدل بازنمایی‌های ویژۀ وظیفه یاد بگیرد.
\item \textbf{ساختار مبتنی بر گراف}: نمایش روابط آزمایشگاه-علائم حیاتی به‌عنوان گراف (مثلاً یال کراتینین-\lr{BUN} برای عملکرد کلیه) و استفاده از شبکه‌های عصبی گرافی برای کدگذاری. این ساختار حوزۀ بالینی را به‌عنوان سوگیری القایی معرفی می‌کند.
\item \textbf{تجمیع زمانی چندمقیاس}: به‌جای تجمیع تکی \pnum{24} ساعته، ویژگی‌ها در پنجره‌های \pnum{4}، \pnum{8}، \pnum{12}، و \pnum{24} ساعت محاسبه شوند تا بازنمایی زمانی غنی‌تر روندها را بگیرد.
\item \textbf{توکن‌سازی ترکیبی}: ترکیب ویژگی‌های جدولی خام با بازنمایی‌های توکن تا مدل به هر دو قالب دسترسی داشته باشد. معماری ترکیبی \lr{ETHOS} و \lr{Random Forest} این راهبرد را تا حدی پیاده می‌کند و به \lr{AUROC} برابر \pnum{0.717} می‌رسد.
\end{enumerate}

%==============================================================================
\section{مقایسۀ بین‌حوزه‌ای: یادگیری کم‌نمونه در انواع داده}
%==============================================================================

برای روشن‌تر شدن جایگاه نتایج، عملکرد یادگیری کم‌نمونه در چند نوع داده با هم مقایسه می‌شود. این مقایسه نشان می‌دهد چرا کم‌نمونه جدولی بالینی از نظر بازنمایی و تعمیم با چالش‌های خاصی روبه‌روست.

\subsection{حوزۀ بینش}

روی \lr{miniImageNet} (\pnum{5}راه، \pnum{5}شات)، شبکه‌های نمونه‌اولیه حدود \pnum{68}\% دقت با ستون فقرات \lr{CNN} \pnum{4}لایه و حدود \pnum{74}\% با ستون فقرات \lr{ResNet-12} به‌دست می‌آورند~\cite{snell2017prototypical}. عملکرد قوی ناشی از: (\pnum{1}) ستون فقرات‌های ازپیش‌آموزش‌دیده قوی با ویژگی‌های غنی و قابل‌انتقال؛ (\pnum{2}) ساختار فضایی در تصاویر که شبکه‌های کانولوشنی به‌کارآمد بهره می‌برند؛ (\pnum{3}) مجموعه‌های فرا آموزش بزرگ (\pnum{64} کلاس، $\sim$\pnum{38,000} تصویر) که نمونه‌برداری اپیزود متنوع را ممکن می‌کند.

\subsection{حوزۀ \lr{NLP}}

در طبقه‌بندی متن کم‌نمونه (\lr{FewRel}، \lr{ARSC})، شبکه‌های نمونه‌اولیه \prange{70}{85}\% دقت بسته به وظیفه و ستون فقرات به‌دست می‌آورند~\cite{han2018fewrel}. مدل‌های زبانی ازپیش‌آموزش‌دیده (\lr{BERT}) به‌عنوان ستون فقرات سود چشمگیری می‌دهند، با \lr{BERT}+\lr{Proto} تنظیم‌شده با دقت بالاتر از \mbox{\lr{$85$\%}} روی \lr{FewRel}. موفقیت ناشی از: (\pnum{1}) بازنمایی‌های ازپیش‌آموزش‌دیده قوی از متن عظیم؛ (\pnum{2}) شباهت معنایی که با طبقه‌بندی مبتنی بر فاصله هم‌راستا است.

\subsection{حوزۀ کشف دارو}

در یادگیری کم‌نمونه مولکولی، به‌ویژه در بنچمارک \lr{FS-Mol}، شبکه‌های نمونه‌اولیه برای پیش‌بینی خاصیت‌ها معمولاً \lr{AUROC} حدود \pnum{0.60} تا \pnum{0.70} به‌دست می‌آورند~\cite{stanley2021fs}. این بازه با خط پایۀ منجمد \texttt{\lr{Proposed\_FewShot}} در این پژوهش (\pnum{0.611}) قابل مقایسه است. چنین شباهتی نشان می‌دهد پیش‌بینی خاصیت مولکولی و امکان‌پذیری بالینی هر دو با ورودی‌های ناهمگن، مرزهای تصمیم غیرخطی و محدودیت دادۀ پیش‌آموزش روبه‌رو هستند.

\subsection{حوزۀ جدولی}

کار اخیر روی یادگیری کم‌نمونه جدولی~\cite{hegselmann2023tabllm} با بردارهای نهفتۀ مدل زبانی (\lr{TabLLM}) نتایج متغیر گزارش می‌کند: رقابتی با خطوط پایه در برخی مجموعه‌ها اما ضعیف‌تر در دیگران. یافته کلی آن است که کم‌نمونه جدولی همچنان چالش باز است، سازگار با نتایج ما.

\subsection{ترکیب}

جدول~\ref{tab:cross_domain} مقایسۀ بین‌حوزه‌ای را خلاصه می‌کند.

\begin{table}[htbp]
\centering
\caption{عملکرد یادگیری کم‌نمونه در حوزه‌های مختلف}
\footnotesize
\setlength{\tabcolsep}{2pt}
\renewcommand{\arraystretch}{1.18}
\begin{tabularx}{\linewidth}{@{}R{1.55cm}C{2.7cm}C{2.2cm}C{2.25cm}C{2.35cm}@{}}
\toprule
\textbf{حوزه} & \textbf{بنچمارک} & \textbf{\lr{Proto-Net}} & \textbf{بهترین} & \shortstack[c]{\textbf{ازپیش}\\\textbf{آموزش؟}} \\
\midrule
بینش & \shortstack[c]{\lr{miniImageNet}\\\lr{5-way, 5-shot}} & \shortstack[c]{\lr{68\%}\\دقت} & \lr{80\%+} & \shortstack[c]{بله\\(\lr{ImageNet})} \\
\lr{NLP} & \shortstack[c]{\lr{FewRel}\\\lr{5-way, 5-shot}} & \shortstack[c]{\lr{72\%}\\دقت} & \lr{85\%+} & \shortstack[c]{بله\\(\lr{BERT})} \\
مولکولی & \lr{FS-Mol} & \shortstack[c]{\cnum{0.65}\\\lr{AUROC}} & \cnum{0.72} & جزئی \\
جدولی (ما) & \shortstack[c]{\lr{MIMIC\text{-}IV}\\\lr{2-way, 5-shot}} & \shortstack[c]{\cnum{0.611}\\\lr{AUROC}} & \shortstack[c]{\textbf{\cnum{0.746}} انباشت\\ \cnum{0.748} سقف جدولی\\ \cnum{0.737} \lr{V5}} & جزئی \\
\bottomrule
\end{tabularx}
\label{tab:cross_domain}
\end{table}

الگوی کلی روشن است: عملکرد کم‌نمونه زمانی بهتر می‌شود که بازنمایی‌های ازپیش‌آموزش‌دیده‌ی قوی وجود داشته باشد و ساختار داده با سوگیری‌های القایی مدل هم‌راستا باشد. دادۀ جدولی بالینی معمولاً هر دو مزیت را کمتر دارد و همین موضوع بخشی از شکاف عملکرد را توضیح می‌دهد. بنابراین برای کاهش این شکاف، یا باید پیش‌آموزش مقیاس‌بزرگ روی دادۀ جدولی بالینی انجام شود، مشابه نقشی که \lr{ImageNet} در بینایی ماشین داشته است، یا باید رویکردهای فرا یادگیری تازه‌ای طراحی شوند که به‌طور خاص با ویژگی‌های جدولی ناهمگن سازگار باشند.

%==============================================================================
\section{پیامدهای بالینی و تحلیل استقرار}
%==============================================================================

\subsection{تأثیر بر گردش‌کار بالینی}

مدل پیش‌بینی امکان‌پذیری درمان می‌تواند در چند نقطۀ تصمیم در گردش‌کار \lr{ICU} یکپارچه شود:

\begin{enumerate}
\item \textbf{ارزیابی \pnum{24}ساعته}: پس از \pnum{24} ساعت اول پذیرش \lr{ICU}، مدل امتیاز امکان‌پذیری بر اساس آزمایش‌ها و علائم حیاتی موجود تولید می‌کند. بیماران با امتیاز بالا (بیش از \pnum{0.7}) برای برنامه‌ریزی زودهنگام ترخیص علامت‌گذاری می‌شوند.
\item \textbf{ارزیابی مجدد روزانه}: با دسترسی داده جدید، مدل پیش‌بینی را به‌روز می‌کند. روند امتیاز امکان‌پذیری (صعودی در برابر نزولی) اطلاعات اضافی دربارۀ مسیر بیمار می‌دهد.
\item \textbf{گردش‌های چندرشته‌ای}: امتیاز امکان‌پذیری همراه دادۀ بالینی دیگر در گردش‌های صبح ارائه می‌شود و تصمیم‌های مشترک برنامه‌ریزی ترخیص را پشتیبانی می‌کند.
\item \textbf{مدیریت تخت}: پیش‌بینی‌های تجمیعی امکان‌پذیری می‌توانند در برنامه‌ریزی ظرفیت \lr{ICU} و برآورد دسترسی به تخت به کار روند.
\end{enumerate}

\subsection{اعتماد و پذیرش بالینی}

برای پذیرش بالینی، چند عامل باید پوشش داده شوند:

\begin{enumerate}
\item \textbf{تفسیرپذیری}: بالینیان باید بفهمند \emph{چرا} مدل امتیاز خاصی پیش‌بینی می‌کند. توضیحات \lr{SHAP} (نشان دادن کراتینین بالا و پلاکت پایین پیش‌بینی «غیرممکن» را هدایت می‌کنند) با استدلال بالینی هم‌راستا است و اعتماد می‌سازد. تحلیل اهمیت ویژگی ما تأیید می‌کند ویژگی‌های پیش‌بینی‌کنندۀ برتر (کراتینین، گلوکز، هموگلوبین) با تعیین‌کننده‌های شناخته‌شدۀ ترخیص هم‌خوان‌اند.

\item \textbf{کالیبراسیون}: مدلی که می‌گوید «\pnum{70}\% احتمال ترخیص امکان‌پذیر» باید تقریباً در \pnum{70}\% مواقع درست باشد. تحلیل کالیبراسیون ما نشان می‌دهد \lr{RF} نسبتاً خوب کالیبره است (\lr{Brier} \pnum{0.21})، در حالی که \lr{ETHOS} گرایش به بیش‌اعتمادی دارد. بازکالیبراسیون (\lr{Platt scaling} یا رگرسیون ایزوتونیک) پیش از استقرار لازم است.

\item \textbf{خستگی هشدار}: اگر مدل برای هر بیمار هشدار تولید کند، بالینیان آن‌ها را نادیده می‌گیرند. انتخاب آستانه باید حساسیت (گرفتن بیماران غیرممکن) و ویژگی (اجتناب از هشدار کاذب) را متعادل کند. \lr{RF} با آستانۀ تنظیم‌شده \pnum{0.54} تعادل عمل‌گرایانه‌ای می‌دهد، اما تنظیم ویژۀ سایت لازم است.

\item \textbf{تخریب تدریجی}: مدل باید دادۀ گم‌شده، پروفایل‌های غیرمعمول، و ورودی‌های خارج از توزیع را بدون پیش‌بینی گمراه‌کننده مدیریت کند. امتناع مبتنی بر اطمینان (امتناع از پیش‌بینی وقتی عدم‌قطعیت مدل بالاست) رویکرد امیدبخشی است.
\end{enumerate}

\subsection{معماری استقرار فدرال}

برای استقرار چندمؤسسه‌ای، معماری زیر را تصور می‌کنیم:

\begin{enumerate}
\item هر مؤسسۀ شرکت‌کننده گرۀ محلی \lr{ETHOS} را اجرا می‌کند که روی دادۀ خود آموزش دیده است.
\item سرور تجمیع مرکزی (ممکن است توسط طرف سوم بی‌طرف میزبانی شود) دورهای \lr{FedAvg} را هفتگی یا ماهانه هماهنگ می‌کند و مدل سراسری را به‌روز می‌کند.
\item مدل‌های محلی می‌توانند پس از دریافت به‌روزرسانی سراسری از راه تنظیم دقیق روی دادۀ ویژۀ سایت شخصی‌سازی شوند و تعمیم سراسری را با سازگاری محلی متعادل کنند.
\item نسخه‌گذاری مدل و پایش انحراف کمک می‌کند مدل استقراریافته با تکامل جمعیت بیماران و شیوه‌های بالینی کالیبره بماند.
\end{enumerate}

شکاف محدود فدرال، حدود \pnum{0.005} واحد \lr{AUROC} در آزمایش گزارش‌شده، این مدل استقرار را پشتیبانی می‌کند: مؤسسات می‌توانند در بهبود مشارکتی مدل شرکت کنند بدون آنکه در این شبیه‌سازی افت بزرگی در عملکرد دیده شود.

\subsection{مسیر نظارتی}

استقرار بالینی نیازمند پیمایش چشم‌انداز نظارتی است:
\begin{itemize}
\item \textbf{\lr{FDA} (\lr{US})}: مدل احتمالاً به‌عنوان نرم‌افزار دستگاه پزشکی ردۀ \lr{II} (\lr{SaMD}) طبق برنامه نوآوری سلامت دیجیتال \lr{FDA} طبقه‌بندی می‌شود. مسیر \lr{510(k)} در صورت وجود دستگاه پیشین ممکن است مناسب باشد.
\item \textbf{علامت \lr{CE} (\lr{EU})}: طبق آیین‌نامه دستگاه پزشکی \lr{EU} (\lr{MDR} \pnum{2017}/\pnum{745})، نرم‌افزار پشتیبانی تصمیم بالینی با تأثیر مستقیم بر بیمار نیازمند ارزیابی انطباق است.
\item \textbf{بررسی مؤسسه‌ای}: اعتبارسنجی آینده‌نگر نیازمند تأیید \lr{IRB} در هر مؤسسۀ شرکت‌کننده است.
\end{itemize}

کار ما شواهد بنیادی (عملکرد مدل، کالیبراسیون، تحلیل انصاف) را که ارسال نظارتی را پشتیبانی می‌کند فراهم می‌کند. با این حال، اعتبارسنجی آینده‌نگر روی کوهورت‌های مستقل پیش‌نیاز هر مسیر نظارتی است.

%==============================================================================
\section{نوآوری و مشارکت‌های اصلی}
%==============================================================================

جنبه‌های نوآورانۀ این پایان‌نامه را که آن را از کار پیشین متمایز می‌کند صریح مرزبندی می‌کنیم:

\begin{enumerate}
\item \textbf{بازنمایی \lr{Symptom Token} برای دادۀ جدولی بالینی.} با وجود کار پیشین روی توکن‌سازی یادداشت‌های بالینی یا سری زمانی، در این پژوهش قالب توکن ساخت‌یافته‌ای برای آزمایش‌ها و علائم حیاتی تجمیع‌شده معرفی می‌شود. هر توکن چهار بعد را کدگذاری می‌کند: هویت علامت، شدت نرمال‌شده، جابه‌جایی زمانی و نوع داده. این بازنمایی، پردازش مبتنی بر ترنسفورمر را برای دادۀ ذاتاً جدولی ممکن می‌کند.

\item \textbf{بررسی کم‌نمونه‌ی نمونه‌اولیه برای پیش‌بینی امکان‌پذیری درمان روی \lr{MIMIC\text{-}IV}.} در چارچوب مرور ادبیات این پایان‌نامه، کاربرد کم‌نمونه بیشتر در تصویربرداری پزشکی (\lr{X-ray}، پوست‌شناسی) و کشف دارو دیده می‌شود و کمتر برای پیش‌بینی ترخیص/امکان‌پذیری درمان \lr{ICU} جدولی ارزیابی شده است. این پژوهش این شکاف را به‌صورت تجربی بررسی می‌کند و هم ظرفیت و هم محدودیت‌های فرا یادگیری اپیزودیک برای دادۀ بالینی ساخت‌یافته را نشان می‌دهد.

\item \textbf{چارچوب یکپارچۀ کم‌نمونه + فدرال + ترنسفورمر.} کارهای پیشین معمولاً شبکه‌های نمونه‌اولیه، میانگین‌گیری فدرال و کدگذاری ترنسفورمر را هم‌زمان برای یک وظیفۀ پیش‌بینی بالینی ترکیب نکرده‌اند. این یکپارچه‌سازی سه چالش را هم‌زمان پوشش می‌دهد: کمبود داده، حفظ حریم خصوصی و بازنمایی ناهمگن داده.

\item \textbf{گزارش صادقانه با حل سازنده.} پیکربندی خط پایه، با \lr{AUROC} حدود \pnum{0.619} در جدول کمپین \lr{V}، ابتدا از خطوط پایۀ کلاسیک عقب می‌ماند. به‌جای کنار گذاشتن این یافته، \emph{چرایی} شکاف تحلیل می‌شود: ناسازگاری اپیزودی، گلوگاه بازنمایی و مقیاس داده. سپس شواهد حذف تدریجی و بهبودهای مرحله‌ای مستند می‌گردد؛ در نهایت اجرای انباشتی \textbf{منجمد} با \lr{AUROC} \textbf{\pnum{0.746}} در برابر \lr{SOTA} جدولی \textbf{\pnum{0.748}} قرار می‌گیرد.

\item \textbf{فرا یادگیر انباشتی و آموزش پیشرفته.} در کمپین \lr{V}، انباشت \lr{V4} با ترکیب احتمال‌های کم‌نمونه، پیش‌بینی‌های \lr{RF} و بردارهای نهفتۀ رمزگذار به \lr{AUROC} برابر \pnum{0.732} می‌رسد. اجرای انباشتی منجمد مخزن \texttt{\lr{exp\_fewshot\_best}} با مقدار \textbf{\pnum{0.746}} ادعای کمی اصلی برای ترکیب کم‌نمونه + کلاسیک در این نسخه است.

\item \textbf{توجه درگاهی به ویژگی و شبکه‌های نمونه‌اولیه \lr{Cross-Attention}.} معماری \lr{V5} انتخاب ویژگی تطبیقی به‌ازای نمونه (\lr{GFA}) و شباهت وابسته به زمینه (نمونه‌های اولیه \lr{Cross-Attention}) را معرفی می‌کند، با بهترین بذر \lr{AUROC} \pnum{0.737} در کمپین توسعه. \lr{Ensemble} ناهمگن \lr{V6} عدم‌قطعیت به‌سبک مطابق را اکتشاف می‌کند.

\item \textbf{سقف تجربی و بنچمارک جامع.} \lr{AUROC} حدود \pnum{0.76} را به‌عنوان سقف تجربی امکان‌پذیری درمان روی این کوهورت معرفی می‌کنیم و نقطۀ مرجعی برای روش‌های آینده فراهم می‌سازیم. ارزیابی ما بیش از \pnum{15} مدل، \pnum{7} مقدار $K$-شات، \pnum{5} گرۀ بیماری، اعتبارسنجی متقاطع \pnum{5}‌لایه، و چند بذر تصادفی را پوشش می‌دهد.
\end{enumerate}

%==============================================================================
\section{جدول مقایسه: کار ما در برابر پیشین}
%==============================================================================

جدول~\ref{tab:comparison_prior} جایگاه این پژوهش را نسبت به رویکردهای پیشین نشان می‌دهد. تفاوت‌های اصلی عبارت‌اند از: ترکیب کم‌نمونه، یادگیری فدرال و بازنمایی مبتنی بر توکن؛ تمرکز بر دادۀ جدولی بالینی به‌جای تصویر؛ و گزارش شفاف نتیجۀ منفی اولیه همراه با مسیر بهبود مرحله‌ای تا نزدیک شدن به سقف جدولی.

\begin{table}[htbp]
\centering
\caption{مقایسه با کار پیشین}
\footnotesize
\setlength{\tabcolsep}{1.8pt}
\renewcommand{\arraystretch}{1.12}
\begin{tabularx}{\linewidth}{@{}R{1.5cm}R{3.0cm}R{3.1cm}Y@{}}
\toprule
\textbf{جنبه} & \textbf{پیشین (معمول)} & \textbf{کار ما} & \textbf{یادداشت} \\
\midrule
داده & \lr{EHR} متمرکز و تجمیع‌شده & آموزش فدرال روی \pnum{5} گرۀ بیماری & حریم خصوصی بهتر حفظ می‌شود، هرچند استقرار واقعی هنوز به اعتبارسنجی عملیاتی نیاز دارد \\
چیدمان & نظارت استاندارد روی کل داده & آموزش اپیزودیک کم‌نمونه با امکان انباشت کلاسیک & برای سناریوهای کم‌داده و توزیع‌شده مناسب‌تر از یک فرایند صرفاً نظارت‌شده است \\
بازنمایی & ویژگی خام، \lr{RNN} یا ترنسفورمر عمومی & \lr{Symptom Tokens} + رمزگذار \lr{ETHOS} & بازنمایی توکنی برای آزمایشگاه و علائم حیاتیِ تجمیع‌شده طراحی شده است \\
وظیفه & مرگ‌ومیر، \lr{LOS}، بستری مجدد & امکان‌پذیری درمان/ترخیص & مسئله به تصمیم‌سازی ترخیص زودهنگام نزدیک‌تر است \\
بهترین کلاسیک & \lr{XGBoost}، \lr{RF} یا \lr{LSTM} روی داده‌ی بزرگ & سقف جدولی \lr{SOTA}/\lr{RF} در بازۀ \prange{0.748}{0.760} & این بازه سقف تجربی کوهورت حاضر را نشان می‌دهد \\
بهترین عصبی & معمولاً در داده‌ی بزرگ رقابتی یا برتر & \lr{Stacked BEST} با \textbf{\pnum{0.746}}؛ \lr{V5}/\lr{V4} به‌عنوان نقاط عطف تاریخی & انباشت منجمد فقط \textbf{\pnum{0.002}} پایین‌تر از برآورد نقطه‌ای \lr{SOTA} است \\
\bottomrule
\end{tabularx}
\label{tab:comparison_prior}
\end{table}

%==============================================================================
\section{ملاحظات اخلاقی و انصاف}
%==============================================================================

استقرار مدل‌های پیش‌بینی بالینی نگرانی‌های اخلاقی برمی‌انگیزد. \textbf{سوگیری}: کوهورت ما ممکن است برخی جمعیت‌شناسی‌ها را کم‌نمایندگی کند؛ عملکرد طبقه‌بندی‌شده سنی (فصل~\ref{ch:data_analysis}) برای ارزیابی انصاف گزارش می‌کنیم. \textbf{شفافیت}: توضیحات \lr{SHAP} و اهمیت ویژگی ارائه می‌دهیم؛ بالینیان باید استدلال مدل را بفهمند. \textbf{پاسخگویی}: مدل‌ها جایگزین قضاوت بالینی نیستند؛ ابزار پشتیبانی تصمیم‌اند. \textbf{حریم خصوصی}: یادگیری فدرال محلیت داده را حفظ می‌کند؛ دادۀ خام بیمار منتقل نمی‌شود. برای استقرار تولید توصیه می‌کنیم: (\pnum{1}) ممیزی انصاف در سن، جنس، و تشخیص؛ (\pnum{2}) پایش مداوم برای جابجایی توزیع؛ (\pnum{3}) انسان-در-حلقه برای تصمیم‌های با ریسک بالا؛ (\pnum{4}) مستندسازی شفاف محدودیت‌ها.

%==============================================================================
\section{چک‌لیست بازتولیدپذیری}
%==============================================================================

به بهترین شیوه‌های بازتولیدپذیری پایبندیم: (\pnum{1}) بذر تصادفی ثابت (\pnum{42}) در همه‌ی آزمایش‌ها؛ (\pnum{2}) تقسیم‌های داده ذخیره و به‌اشتراک گذاشته شده؛ (\pnum{3}) شبکه‌های ابرپارامتر در پیوست~\ref{app:hyperparams} مستند شده‌اند؛ (\pnum{4}) ساختار کد چیدمان استاندارد پروژه را دنبال می‌کند؛ (\pnum{5}) دسترسی به \lr{MIMIC\text{-}IV} نیازمند احراز هویت \lr{PhysioNet} است. پژوهشگران می‌توانند با دسترسی \lr{MIMIC\text{-}IV}، نصب وابستگی‌ها از \texttt{\lr{requirements.txt}}، و اجرای اسکریپت‌های آزمایش با پیکربندی مستند، نتایج را بازتولید کنند.

%==============================================================================
\section{درس‌های آزمایش‌های \lr{V5} و \lr{V6}}
%==============================================================================

آزمایش‌های \lr{V5} و \lr{V6} چند نکته‌ی مهم برای یادگیری کم‌نمونه روی دادۀ جدولی نشان می‌دهند:

\textbf{توجه درگاهی به ویژگی برای دادۀ بالینی مفید است.} ماژول \lr{GFA} در \lr{V5}، با الهام از \lr{TabNet}، برای هر نمونه ماسک‌های اهمیت ویژگی می‌آموزد. بهترین بذر \lr{V5} در کمپین \lr{V} به \lr{AUROC} \pnum{0.737} می‌رسد و سقف بالقوۀ مسیر کم‌نمونه را نشان می‌دهد؛ البته اجرای منجمد بهترین زیرمدل کم‌نمونه با \lr{GFA} و \lr{Cross-Attention} با \lr{AUROC} \textbf{\pnum{0.670}} مرجع رسمی مسیر فقط کم‌نمونه در این نسخه است. واریانس بالاتر میان بذرها (\lr{AUROC} \cnum{0.703}-\cnum{0.737} در برابر \cnum{0.696}-\cnum{0.713} برای \lr{V4}) نشان می‌دهد این معماری به مقداردهی اولیه حساس‌تر است.

\textbf{نمونه‌های اولیه با \lr{Cross-Attention} از نظر نظری انعطاف‌پذیرترند.} سر \lr{Cross-Attention} به‌جای تکیه بر فاصلۀ کسینوسی ثابت، شباهتی وابسته به زمینه می‌آموزد؛ رویکردی که برای دادۀ بالینی ناهمگن و مرزهای چندوجهی مناسب‌تر است. با این حال، سود آن روی مجموعۀ فعلی با \lr{$n=2{,}000$} محدود بود. این نتیجه نشان می‌دهد \lr{Cross-Attention} احتمالاً برای آشکار کردن ظرفیت کامل خود به مجموعه‌های بزرگ‌تر نیاز دارد.

\textbf{همگنی معماری در مقیاس کوچک بر تنوع غلبه می‌کند.} \lr{Ensemble} ناهمگن \lr{V6} (\pnum{0.725}) از \lr{Ensemble}های هم‌معماری \lr{V4}، \lr{V5} (\pnum{0.732}، \pnum{0.730}) در کمپین \lr{V} عقب می‌ماند. معماری \lr{MLP-Mixer} (\lr{AUROC} \pnum{0.642}) \lr{Ensemble} را پایین می‌کشد چون فاقد خودتوجه است؛ تأیید می‌کند برای دادۀ جدولی با تعامل پیچیده بین‌ویژگی، مکانیزم‌های توجه اختیاری نیستند.

\textbf{انباشت ساده در مقیاس کوچک اعتبارسنجی از انباشت پیچیده بهتر است.} فرا یادگیر \lr{XGBoost} \lr{V6} (\lr{AUROC} \pnum{0.689}) روی مجموعۀ اعتبارسنجی \pnum{320}نمونه‌ای بیش‌برازش می‌کند، در حالی که رگرسیون لجستیک کالیبره‌شده ایزوتونیک \lr{V4} (\pnum{0.732}) در آن کمپین خوب تعمیم می‌یابد. پیکربندی انباشتی منجمد \texttt{\lr{exp\_fewshot\_best}} باید برای عملکرد انباشت سطح ارسال استناد شود.

%==============================================================================
\section{کارهای آینده}
%==============================================================================

\begin{enumerate}
\item \textbf{پیش‌آموزش بزرگ‌تر}: آموزش \lr{ETHOS} روی کل کوهورت \lr{MIMIC\text{-}IV} (بیش از \pnum{50,000} اقامت در \lr{ICU}) یا دادۀ \lr{EHR} خارجی پیش از تنظیم دقیق کم‌نمونه. اوج بذرهای \lr{V5} در کمپین \lr{V} به \lr{AUROC} \pnum{0.737} روی \pnum{1,280} نمونۀ آموزش رسید؛ مجموعه‌های بزرگ‌تر ممکن است عملکرد فقط کم‌نمونه را فراتر از برآورد نقطه‌ای منجمد \pnum{0.670} تثبیت و بهبود دهند.
\item \textbf{مقیاس‌دهی \lr{GFA} و \lr{Cross-Attention}}: ارزیابی معماری \lr{V5} روی مجموعه‌های بالینی بزرگ‌تر جایی که نمونه‌های اولیه \lr{Cross-Attention} سود بیشتری دهند.
\item \textbf{انباشت ترکیبی \lr{GFA}}: ترکیب رمزگذار \lr{GFA} \lr{V5} با انباشت کالیبره‌شده \lr{V4} برای بهترین هر دو رویکرد.
\item \textbf{ارزیابی فدرال بین‌مؤسسه‌ای}: ارزیابی روی دادۀ چند بیمارستان برای ارزیابی تعمیم، به‌ویژه مکانیزم سازگاری \lr{Proto-MAML}.
\item \textbf{پیش‌بینی مطابق در استقرار}: گسترش پیش‌بینی مطابق \lr{V6} برای برآورد عدم‌قطعیت سطح بیمار در سامانۀ پشتیبانی تصمیم بالینی.
\item \textbf{مدل‌سازی زمانی}: گنجاندن پویایی سری زمانی (روند بهبود در برابر وخامت) برای کاهش خطا در موارد مرزی.
\item \textbf{مقیاس‌دهی چندوظیفه}: گسترش به وظایف بیشتر (مثلاً شروع سپسیس، \lr{AKI}) با رمزگذارهای بزرگ‌تر.
\item \textbf{اعتبارسنجی آینده‌نگر}: استقرار در مطالعۀ آزمایشی با بازخورد بالینی و ردیابی پیامد.
\item \textbf{\lr{MLP-Mixer} تقویت‌شده با توجه}: بررسی معماری‌های ترکیبی توجه-میکسر که کارایی مخلوط‌کردن توکن را با بیانگری خودتوجه ترکیب می‌کنند.
\end{enumerate}

جدول~\ref{tab:future_roadmap} نقشه‌راه مرحله‌ای کار آینده را نشان می‌دهد. در فاز \pnum{1}، تمرکز بر مقیاس‌دهی به کل \lr{MIMIC\text{-}IV} و اعتبارسنجی خارجی روی \lr{eICU} است. در فاز \pnum{2}، آزمایش فدرال چندسایتی، فرا یادگیرهای جایگزین مانند \lr{MAML} و مقیاس‌دهی ترکیبی بررسی می‌شوند. فاز \pnum{3} به استقرار آینده‌نگر، مسیرهای نظارتی و دریافت بازخورد بالینی اختصاص دارد. هر فاز بر خروجی فاز پیشین بنا می‌شود و پایان‌نامۀ حاضر پایۀ فنی فاز نخست را فراهم می‌کند.

\begin{table}[htbp]
\centering
\caption{نقشه‌راه کار آینده}
\footnotesize
\setlength{\tabcolsep}{2pt}
\renewcommand{\arraystretch}{1.14}
\begin{tabularx}{\linewidth}{@{}C{1.2cm}R{2.1cm}Y@{}}
\toprule
\textbf{فاز} & \textbf{زمان‌بندی} & \textbf{تمرکز} \\
\midrule
\pnum{1} & \pnum{1} سال & \lr{MIMIC\text{-}IV} کامل، اعتبارسنجی \lr{eICU}، پیش‌آموزش بزرگ‌تر \\
\pnum{2} & \pnum{2} سال & فدرال چندسایتی، \lr{MAML}، مقیاس‌دهی ترکیبی \\
\pnum{3} & \pnum{3} سال & استقرار آینده‌نگر، نظارتی، بازخورد بالینی \\
\bottomrule
\end{tabularx}
\label{tab:future_roadmap}
\end{table}

نقشه‌راه مقیاس‌دهی (فاز \pnum{1})، کاوش فدرال و معماری (فاز \pnum{2})، و استقرار (فاز \pnum{3}) را اولویت‌بندی می‌کند. هر فاز به منابع و مشارکت‌های اختصاصی نیاز دارد؛ پایان‌نامۀ فعلی بنیاد فنی را فراهم می‌کند.

%==============================================================================
\section{خلاصۀ نکات کلیدی و درس‌های آموخته‌شده}
%==============================================================================

نکات اصلی را برای پژوهشگران و عملگران استخراج می‌کنیم:

\begin{enumerate}
\item \textbf{مدل‌های کلاسیک روی دادۀ جدولی متوسط سقف محکمی می‌سازند.} \lr{Random Forest} به \lr{AUROC} حدود \pnum{0.76} می‌رسد؛ بنابراین برای ارزیابی منصفانۀ روش‌های جدید، داشتن خط پایۀ کلاسیک قوی ضروری است.
\item \textbf{فرایندهای کم‌نمونه با مهندسی سیستماتیک می‌توانند به \lr{SOTA} جدولی نزدیک شوند.} شکاف میان کم‌نمونه ساده و مدل‌های جدولی قوی در ابتدا زیاد است، اما انباشت، تقطیر، سازگاری و اصلاح معماری آن را کاهش می‌دهند. اجرای انباشتی منجمد با \textbf{\pnum{0.746}} در برابر \lr{SOTA} \textbf{\pnum{0.748}} مبنای کمی ادعای نزدیک شدن به سقف است؛ در حالی که بذرهای کمپین \lr{V} مانند \pnum{0.737} رفتار تاریخی بهترین اجرای تکی را نشان می‌دهند.
\item \textbf{یادگیری فدرال عملی است} با شکاف محدود حدود \pnum{0.005} واحد \lr{AUROC} در آزمایش گزارش‌شده.
\item \textbf{فرا یادگیرهای انباشتی بهترین راهبرد آمیزش} در مقیاس کوچک اعتبارسنجی ($n=320$) هستند. رگرسیون لجستیک کالیبره‌شده از انباشت \lr{XGBoost} بهتر عمل می‌کند.
\item \textbf{تنوع معماری در مقیاس کوچک کمک نمی‌کند.} \lr{Ensemble}های همگن (همان معماری، بذرهای متفاوت) از \lr{Ensemble}های ناهمگن بهترند. مکانیزم‌های توجه برای یادگیری کم‌نمونه جدولی ضروری‌اند.
\item \textbf{تفسیرپذیری در بهداشت مهم است.} \lr{SHAP} و اهمیت ویژگی برای اعتماد بالینی ضروری‌اند.
\item \textbf{مهندسی ویژگی ظرفیت مدل عصبی را تقویت می‌کند.} گسترش از \pnum{22} به \pnum{132} ویژگی سیگنال غنی‌تری برای رمزگذار فراهم می‌کند.
\end{enumerate}

پرسش‌های باز برای کار آینده عبارت‌اند از: (\pnum{1}) بررسی اینکه پیش‌آموزش بزرگ‌تر تا چه اندازه عملکرد رمزگذار \lr{GFA} را بهبود می‌دهد؛ (\pnum{2}) اعتبارسنجی خارجی روی دادۀ چندمرکز (\lr{eICU})؛ (\pnum{3}) ترکیب رمزگذار \lr{GFA} \lr{V5} با انباشت \lr{V4}؛ (\pnum{4}) مقیاس‌دهی پیش‌بینی مطابق برای سنجش عدم‌قطعیت آمادۀ استقرار؛ و (\pnum{5}) اعتبارسنجی بالینی آینده‌نگر.

%==============================================================================
\section{خلاصۀ محدودیت‌ها}
%==============================================================================

محدودیت‌های زیر را می‌پذیریم: (\pnum{1}) \textbf{تک مرکز}: \lr{MIMIC\text{-}IV} از یک بیمارستان (\lr{BIDMC}) است؛ نتایج ممکن است تعمیم نیابند. (\pnum{2}) \textbf{اندازۀ نمونه}: \pnum{2,000} بیمار متوسط است؛ \lr{MIMIC\text{-}IV} کامل (\pnum{250}هزار+ اقامت) می‌تواند نتیجه‌گیری متفاوتی بدهد. (\pnum{3}) \textbf{تعریف برچسب}: امکان‌پذیری دودویی است (مقصد ترخیص + بستری مجدد)؛ تعاریف جایگزین نتایج را عوض می‌کند. (\pnum{4}) \textbf{ویژگی‌ها}: ویژگی‌های پایه \pnum{25} بعد (آزمایشگاه، علائم حیاتی، مهندسی‌شده)، گسترش به \pnum{132} با مهندسی ویژگی؛ یادداشت‌های بالینی، تصویربرداری، و داروها مستثنی‌اند. (\pnum{5}) \textbf{وضوح زمانی}: تجمیع \pnum{24} ساعت اول؛ روند ساعتی یا زیرساعتی می‌تواند سیگنال اضافه کند. (\pnum{6}) \textbf{شبیه‌سازی فدرال}: پارتیشن فدرال را به‌تفکیک گرۀ بیماری شبیه‌سازی می‌کنیم؛ استقرار چندسایتی واقعی ممکن است چالش‌های اضافی (هزینۀ ارتباط، خرابی گره، مشتریان بیزانسی) داشته باشد. (\pnum{7}) \textbf{واریانس بذر \lr{V5}، \lr{V6}}: رمزگذار \lr{GFA} \lr{V5} واریانس بالای به‌ازای بذر (\lr{AUROC} \cnum{0.703}-\cnum{0.737}) دارد و نتیجۀ بهترین بذر (\cnum{0.737}) ممکن است خوش‌بینانه باشد؛ میانگین‌های \lr{Ensemble} (\cnum{0.722}) نماینده‌ترند. با وجود این محدودیت‌ها، کار ما روش‌شناسی دقیق و نقشه‌راه جامع برای پژوهش آینده فراهم می‌کند.

%==============================================================================
\section{جمع‌بندی فصل}
%==============================================================================

این فصل نتایج تجربی را تفسیر کرد: سقف جدولی منجمد با \textbf{\pnum{0.748}}، اجرای انباشتی کم‌نمونه با \textbf{\pnum{0.746}}، شکاف کنترل‌شدۀ \lr{Exp1} میان مدل پیشنهادی کم‌نمونه با \textbf{\pnum{0.611}} و \lr{RF} با \textbf{\pnum{0.762}}، هزینۀ فدرال در حدود \pnum{0.005} \lr{AUROC}، و مسیر توسعه‌ی \lr{V1--V6}. همچنین محدودیت‌های فرا یادگیری اپیزودیک، مقایسه با کارهای پیشین، ملاحظات استقرار، اخلاق و مسیرهای آینده بررسی شد. با وجود نتایج امیدبخش، استقرار بالینی همچنان به اعتبارسنجی آینده‌نگر و ارزیابی در سطح هر مؤسسه نیاز دارد.
